% ID: SEL-D0014-BC-E06
% ESTADO: cotejada. Segunda pasada visual completa y borrador compilado revisados.
% Fuente: Selectividad/matematicas_aplicadas/matematicas_aplicadas_P_B.pdf
% SHA-256: 51601c4f5a6c1dfcd860f621a659f8b9c09f53ff4cd31a1fa1e827dbd8f06159
% Página: 2; bloque C; ejercicio original 6.

Una empresa almacena el mismo número de latas de refresco de cola, naranja
y limón. De las $30\,000$ latas de refresco almacenadas, se sabe que $1\,800$
latas de cola, $2\,400$ de naranja y $3\,000$ de limón caducan en 2021.
\begin{enumerate}[label=\alph*)]
  \item \textbf{(1.5 puntos)} \textquestiondown{}Cuál es la probabilidad de
  que una lata elegida al azar caduque en 2021?
  \item \textbf{(1 punto)} Si se ha elegido al azar una lata que no caduca
  en 2021, \textquestiondown{}cuál es la probabilidad de que sea de cola?
\end{enumerate}
